\section{The kinematic selector}
\label{sec:selector}

Choosing a Lorentzian four-plane inside $\R^{3,3}$ is equivalent to choosing its orthogonal complement, a two-plane, and a two-plane is described by its unit
bivector. The question is which bivectors describe a Lorentzian complement, and what the arena does with them.

\subsection{One constituent}

\begin{theorem}[The selector is a complex structure]
\label{thm:selector}
Let $C=a\wedge b$ be a simple bivector of $\R^{3,3}$ and $N=\langle a,b\rangle$.
\begin{enumerate}[label=(\roman*),leftmargin=2.2em]
\item $C^2=-\det\mathrm{Gram}(a,b)\,\id$.
\item For a nondegenerate $N$ with $C$ normalised: $C^2=-\id$ exactly when $N$ is definite, that is exactly when $N^\perp$ is Lorentzian, of signature $(1,3)$
when $N$ is positive and $(3,1)$ when $N$ is negative; $C^2=+\id$ exactly when $N^\perp$ has signature $(2,2)$; and $C^2=0$ exactly when $N$ is degenerate.
\item The centraliser of $C$ in $\so{3,3}$ is $\so N\oplus\so{N^\perp}$, of dimension seven, so the choices of $N$ of a given type form an eight-dimensional homogeneous
space.
\item The surplus plane is such an $N$, with $C=J=e_2e_3$ and $N^\perp$ the interior four-plane.
\end{enumerate}
\end{theorem}

\begin{proof}
(i) is the Clifford identity $(a\wedge b)^2=(a\cdot b)^2-a^2b^2$. (ii) For an orthonormal pair, $(ab)^2=-a^2b^2$, which is $-1$ when $a^2=b^2$ and $+1$ otherwise, and the
signature of $N^\perp$ is the complement of that of $N$ in $(3,3)$. (iii) holds for coordinate planes by \eqref{eq:bladecomm}, and every nondegenerate plane is an
$SO(3,3)$-image of a coordinate plane of the same signature, by Witt's theorem.
\end{proof}

For coordinate planes $C$ is moreover antisymmetric on $S$ when $C^2=-\id$, an orthogonal complex structure; for tilted planes it is not. Orthogonality refers to the
Euclidean structure of $S$, which singles out the temporal triple, while $C^2=-\id$ is invariant. The distinction will matter.

Spinors give the light cone. The volume element $\omega_S=e_4e_5e_6$ satisfies $\omega_Se_A\omega_S^{-1}=-e_A^{\!\top}$ for every generator, so
$\langle\psi,\varphi\rangle=\psi^{\!\top}\omega_S\varphi$ is a $\mathrm{Spin}(3,3)$-invariant pairing, and it pairs $S^+$ with $S^-$.

\begin{theorem}[The light cone as twisted squares]
\label{thm:cone}
Let $\psi$ range over the chiral half $S^+$.
\begin{enumerate}[label=(\roman*),leftmargin=2.2em]
\item The forms $m_A=\omega_Se_A$ are antisymmetric and satisfy $X^{\!\top}m_A+m_AX=-\omega_S[X,e_A]$ for every bivector $X$; their restrictions to $S^+$ span
$\Lambda^2(S^+)^*$, so $\psi\wedge\varphi\mapsto(\psi^{\!\top}\omega_Se_A\varphi)_A$ is an equivariant isomorphism $\Lambda^2S^+\cong\R^{3,3}$, the Klein correspondence of
\cref{thm:klein}.
\item On $S^+$, $J^2=-\id$, so $S^+\cong\C^2$, and the interior $\so{1,3}$ acts $\C$-linearly, as $\slr{2,\C}$.
\item For a unit simple bivector $C$ with plane $N$ define $v_C(\psi)_A=\psi^{\!\top}\omega_Se_AC\psi$. Then $v_C(\psi)\in N^\perp$ and $\eta(v_C,v_C)=0$ identically,
and the image has dimension three.
\item If $C^2=-\id$, the component of $v_C$ along the timelike direction of $N^\perp$ is definite: $v_C$ fills one nappe of the light cone of $N^\perp$. If
$C^2=+\id$, no component is definite, and $v_C$ fills the null vectors of a $(2,2)$ plane, which has no time orientation.
\item $v_{-C}=-v_C$; on $S^-$ the $J$-twisted square lies in the other nappe; and on $S^+$, $v_J=-j$, where $j_A=\psi^{\!\top}\gamma^0e_A\psi$ is the Dirac current.
\end{enumerate}
\end{theorem}

\begin{figure}[t]
\centering
\begin{tikzpicture}[>=Stealth,font=\small]
\begin{scope}
\fill[tcol!10] (0,0) -- (1.6,2.2) arc (0:180:1.6 and 0.35) -- cycle;
\draw[tcol,thick] (-1.6,2.2) -- (0,0) -- (1.6,2.2);
\draw[tcol,thick] (0,2.2) ellipse (1.6 and 0.35);
\draw[tcol!60,dashed] (-1.6,-2.2) -- (0,0) -- (1.6,-2.2);
\draw[tcol!60,dashed] (0,-2.2) ellipse (1.6 and 0.35);
\draw[->,gray] (0,-2.7) -- (0,2.9) node[above] {$\gamma^0$};
\node[tcol,align=center] at (0,3.6) {$C=J$, $C^2=-\id$: complement $(1,3)$};
\node[anchor=west,align=left] at (1.75,1.7) {$v_J(\psi)$, $\psi\in S^+$:\\ one nappe};
\node[anchor=west,align=left,tcol!70] at (1.75,-1.7) {$S^-$ or $-J$:\\ the other nappe};
\end{scope}
\begin{scope}[xshift=8.2cm]
\draw[scol,thick] (-1.6,-1.6) -- (1.6,1.6); \draw[scol,thick] (-1.6,1.6) -- (1.6,-1.6);
\draw[scol,thick,dashed] (-2.1,-0.7) -- (2.1,0.7); \draw[scol,thick,dashed] (-2.1,0.7) -- (2.1,-0.7);
\fill[scol!10] (0,0) circle (0.25);
\node[scol,align=center] at (0,3.6) {$C^2=+\id$: complement $(2,2)$};
\node[align=center] at (0,-2.6) {null vectors of a $(2,2)$ plane:\\ connected, no past or future};
\end{scope}
\end{tikzpicture}
\caption{The selector at one constituent (\cref{thm:selector,thm:cone}). A unit bivector $C$ with $C^2=-\id$ selects a Lorentzian complement, and the $C$-twisted squares of
chiral spinors fill one nappe of its light cone, the other nappe belonging to the other chirality or to $-C$. A bivector with $C^2=+\id$ selects a $(2,2)$ complement,
whose null set has no time orientation.}
\label{fig:selector}
\end{figure}

The selector of the reduction is therefore not a new field. Choosing the spacetime $N^\perp$ and choosing the complex structure $C$ are one choice, and in the arena that
complex structure is the charge $J$, the same $J$ that is the imaginary unit of the Dirac theory (\cref{thm:i}). Null directions are the squares of Weyl spinors twisted
by $J$, and the time orientation is fixed by the sign of $J$ together with the chirality. The mathematics is the Klein--Penrose correspondence, in which a null vector is
$\psi\bar\psi$ \cite{PenroseRindler84,WardWells}; what is specific to the arena is that $J$ already plays this role. All of this is kinematics: nothing yet says why $J$,
rather than an indefinite or degenerate bivector, is chosen. For degenerate planes $v_C$ is still null with a three-dimensional image, and for the totally null plane
$\langle e_1+e_4,e_2+e_5\rangle$ it lies in the totally null three-plane $\langle e_i+e_{i+3}\rangle$, an eigenspace of the grading root $\cK$.

\subsection{Three constituents}

At three constituents $J$ is the product of the fourteen generators of the Lorentz complement $W_3$, of signature $(8,6)$, and $\Omega$ is the volume element of $E$.

\begin{theorem}[The selector at three constituents]
\label{thm:selector3}
\begin{enumerate}[label=(\roman*),leftmargin=2.2em]
\item $J^2=\Omega^2=-\id$, $[J,\Omega]=0$ and $\Gamma=J\Omega$ with $\Gamma^2=\id$; $\Omega=-J$ on $\Gamma=+1$. Since $\Gamma^2=\id$ at every rung, any coordinate four-plane $E'$
with volume element $\Omega_{E'}$ and complementary volume element $J_{E'}$ has $J_{E'}^2=\Omega_{E'}^2=\prod_{i\in E'}e_i^2$: the volume element of the complement is a complex
structure exactly when $E'$ is Lorentzian, at every $n$.
\item The centraliser of $J$ among the $153$ bivectors has dimension $97=6+91$, so the choices form a space of dimension $56$.
\item With either $\mathrm{Spin}(9,9)$-invariant pairing, the product $\beta_{\rm odd}$ of the nine negative generators, which is antisymmetric, or the product $\beta_{\rm even}$
of the nine positive ones, which is symmetric, $\beta e_AJ$ is symmetric exactly for the fourteen directions of $W_3$: the $J$-twisted square lives on the complement,
not on $E$. The untwisted $\beta_{\rm odd}e_A$ is symmetric for all eighteen directions, so the vector sits in $\mathrm{Sym}^2S^+$, and at random spinors the untwisted
square is not null.
\end{enumerate}
\end{theorem}

The identification of the selector with $J$ survives at three constituents, but the covariant construction of the cone from spinors is a feature of dimension six. At
$n=3$ the cone of $E$ is recovered by the Dirac current $j_\mu=\psi^{\!\top}\gamma^0\gamma_\mu\psi$, which is future causal, being a sum of one-constituent currents, and
null on product states; but it uses the clock $\gamma^0$ and the transpose, that is, the positive inner product of the spinor module.
